\documentclass[preprint,12pt]{elsarticle}

\usepackage{amsmath,amssymb,amsthm}
\usepackage{booktabs}
\usepackage{longtable}
\usepackage{hyperref}
\pdfstringdefDisableCommands{\def\corref#1{}}

\newtheorem{theorem}{Theorem}
\newtheorem{corollary}{Corollary}

\newcommand{\Q}{Q_6}
\newcommand{\M}{\mathcal M}

\journal{Discrete Mathematics}

\begin{document}

\begin{frontmatter}

\title{The missing \texorpdfstring{\(K_{3,3}\)}{K3,3} case for semi-perfect 1-factorizations of the hypercube}

\author[aff1]{Zhu Mindong\corref{cor1}}
\ead{13788998119@163.com}
\ead[url]{https://orcid.org/0009-0006-6900-305X}
\cortext[cor1]{Corresponding author}

\affiliation[aff1]{
  organization={Institute of Marxism, Postdoctoral Research Station of Sociology, East China University of Science and Technology},
  addressline={130 Meilong Road},
  city={Shanghai},
  postcode={200237},
  country={China}
}

\begin{abstract}
Behague proved that for all positive integers \(k,l\), except for the possible
exception \((k,l)=(3,3)\), there is a 1-factorization \(\M\) of \(Q_{k+l}\)
whose auxiliary graph \(G[\M]\) is isomorphic to \(K_{k,l}\). We give an
explicit 1-factorization of \(Q_6\) for which \(G[\M]\cong K_{3,3}\), thereby
settling the remaining case. The construction is supplied as a compact
edge-coloring table, together with a deterministic verification procedure and
machine-readable certificate. The proof reduces to checking that the six
listed color classes are perfect matchings and that the pairwise cycle lengths
have the prescribed form.
\end{abstract}

\begin{keyword}
hypercube \sep 1-factorization \sep Hamilton cycle \sep graph factorization
\sep combinatorial construction
\end{keyword}

\end{frontmatter}

\section{Introduction}

Let \(Q_d\) be the \(d\)-dimensional hypercube. A 1-factorization of \(Q_d\)
is a partition of \(E(Q_d)\) into \(d\) perfect matchings. A classical theme
in graph factorization asks when the union of two factors is a Hamilton cycle.
In its strongest form this leads to perfect 1-factorizations, a subject going
back to Kotzig's work and conjectures on complete graphs
\cite{kotzig,anderson}. Related existence and non-existence questions have
since been studied for complete bipartite graphs and other highly symmetric
families; see, for example, \cite{laufer,bryant_maenhaut_wanless}.

For hypercubes, the perfect condition is too restrictive, and the natural
question becomes which pairs of factors can be made Hamiltonian. Following
Behague \cite{behague}, if
\[
  \M=\{M_1,\dots,M_d\}
\]
is a 1-factorization, define \(G[\M]\) to be the graph with vertex set
\(\M\), where \(M_iM_j\) is an edge exactly when \(M_i\cup M_j\) is a
Hamilton cycle of \(Q_d\). Thus \(G[\M]\) records the Hamiltonian pair
structure of the factorization. This language is closely related to the
``semi-perfect'' and \(k\)-semi-perfect factorizations studied by Gochev and
Gotchev \cite{gochev_gotchev}, Kr\'alovi\v{c} and Kr\'alovi\v{c}
\cite{kralovic_kralovic}, Chitra and Muthusamy \cite{chitra_muthusamy}, and
others. Hamilton decompositions of directed cubes and related products provide
another source of constructions and obstructions \cite{stong}.

Behague proved that for every 1-factorization \(\M\) of a hypercube the graph
\(G[\M]\) is bipartite, and then considered the converse realization problem:
which bipartite graphs arise as \(G[\M]\)? In particular, he proved that for
all positive integers \(k,l\), not both equal to \(3\), there is a
1-factorization \(\M\) of \(Q_{k+l}\) with \(G[\M]\cong K_{k,l}\), and asked
whether the remaining case \(K_{3,3}\) occurs \cite[Theorem 2 and Open Problem
8]{behague}.

This note gives an explicit construction for the missing case. The proof is
entirely finite: Table~\ref{tab:coloring} specifies the factorization, and the
cycle-length calculation displayed in the proof identifies exactly which
pairs of factors are Hamiltonian. The consequence is that Behague's existence
result for complete bipartite auxiliary graphs holds for all positive
\(k,l\), with no exceptional pair.

\section{The construction}

Identify the vertices of \(\Q\) with the integers \(0,\dots,63\), or
equivalently with the binary vectors in \(\mathbb F_2^6\). Direction \(d\)
means flipping bit \(d\), so the edge in direction \(d\) from \(v\) joins
\(v\) to \(v\oplus 2^d\). Every edge has exactly one endpoint of even parity.

Table~\ref{tab:coloring} gives an edge-coloring of \(\Q\). For an even vertex
\(v\), the entry in column \(d\) is the color of the edge
\(\{v,v\oplus 2^d\}\). Let \(M_i\) be the set of edges of color \(i\).

\begin{longtable}{rrrrrrrr}
\caption{The color of the edge in each direction from each even vertex.}
\label{tab:coloring}\\
\toprule
\(v\) & binary & dir 0 & dir 1 & dir 2 & dir 3 & dir 4 & dir 5 \\
\midrule
\endfirsthead
\toprule
\(v\) & binary & dir 0 & dir 1 & dir 2 & dir 3 & dir 4 & dir 5 \\
\midrule
\endhead
0  & 000000 & 4 & 5 & 0 & 1 & 3 & 2 \\
3  & 000011 & 1 & 0 & 3 & 2 & 4 & 5 \\
5  & 000101 & 2 & 0 & 1 & 5 & 4 & 3 \\
6  & 000110 & 2 & 3 & 4 & 1 & 0 & 5 \\
9  & 001001 & 2 & 4 & 0 & 5 & 1 & 3 \\
10 & 001010 & 3 & 4 & 5 & 2 & 1 & 0 \\
12 & 001100 & 3 & 4 & 0 & 5 & 2 & 1 \\
15 & 001111 & 3 & 2 & 0 & 4 & 1 & 5 \\
17 & 010001 & 2 & 0 & 1 & 5 & 3 & 4 \\
18 & 010010 & 3 & 1 & 2 & 4 & 0 & 5 \\
20 & 010100 & 2 & 5 & 0 & 1 & 4 & 3 \\
23 & 010111 & 4 & 3 & 5 & 2 & 1 & 0 \\
24 & 011000 & 2 & 0 & 5 & 4 & 3 & 1 \\
27 & 011011 & 3 & 0 & 4 & 1 & 5 & 2 \\
29 & 011101 & 4 & 5 & 3 & 0 & 1 & 2 \\
30 & 011110 & 0 & 3 & 5 & 1 & 2 & 4 \\
33 & 100001 & 0 & 1 & 5 & 4 & 3 & 2 \\
34 & 100010 & 0 & 5 & 2 & 1 & 4 & 3 \\
36 & 100100 & 4 & 0 & 3 & 2 & 5 & 1 \\
39 & 100111 & 1 & 0 & 2 & 3 & 4 & 5 \\
40 & 101000 & 1 & 3 & 0 & 4 & 2 & 5 \\
43 & 101011 & 4 & 0 & 2 & 3 & 5 & 1 \\
45 & 101101 & 3 & 1 & 5 & 2 & 0 & 4 \\
46 & 101110 & 4 & 5 & 2 & 3 & 1 & 0 \\
48 & 110000 & 0 & 2 & 4 & 3 & 1 & 5 \\
51 & 110011 & 3 & 5 & 1 & 0 & 4 & 2 \\
53 & 110101 & 0 & 3 & 2 & 4 & 1 & 5 \\
54 & 110110 & 5 & 2 & 1 & 0 & 4 & 3 \\
57 & 111001 & 0 & 3 & 5 & 1 & 2 & 4 \\
58 & 111010 & 1 & 4 & 3 & 0 & 5 & 2 \\
60 & 111100 & 3 & 2 & 5 & 1 & 4 & 0 \\
63 & 111111 & 5 & 1 & 4 & 2 & 0 & 3 \\
\bottomrule
\end{longtable}

\begin{theorem}
The color classes \(M_0,\dots,M_5\) defined by Table~\ref{tab:coloring} form
a 1-factorization of \(Q_6\). Moreover,
\[
  G[\M]\cong K_{3,3},
\]
with bipartition \(\{M_0,M_1,M_5\}\cup\{M_2,M_3,M_4\}\).
\end{theorem}

\begin{proof}
Each row of Table~\ref{tab:coloring} is a permutation of
\(\{0,1,2,3,4,5\}\), so each even vertex is incident with exactly one edge of
each color. A direct check of the table also shows that each odd vertex is
incident with exactly one edge of each color. Hence each \(M_i\) is a perfect
matching, and the six \(M_i\)'s partition the edge set of \(Q_6\).

For each unordered pair \(\{i,j\}\), the graph \(M_i\cup M_j\) is 2-regular
on all 64 vertices, so it is a disjoint union of cycles. The cycle lengths are
as follows:
\[
\begin{array}{c|l}
01 & 40,12,8,4\\
02 & 64\\
03 & 64\\
04 & 64\\
05 & 26,14,12,12\\
12 & 64\\
13 & 64\\
14 & 64\\
15 & 42,22\\
23 & 40,14,6,4\\
24 & 50,6,4,4\\
25 & 64\\
34 & 52,12\\
35 & 64\\
45 & 64
\end{array}
\]
Thus \(M_i\cup M_j\) is Hamiltonian exactly for the nine pairs
\[
  \{0,2\},\{0,3\},\{0,4\},
  \{1,2\},\{1,3\},\{1,4\},
  \{5,2\},\{5,3\},\{5,4\}.
\]
These are precisely the edges of the complete bipartite graph with parts
\(\{0,1,5\}\) and \(\{2,3,4\}\). Therefore \(G[\M]\cong K_{3,3}\).
\end{proof}

\begin{corollary}
For all positive integers \(k,l\), there is a 1-factorization \(\M\) of
\(Q_{k+l}\) such that \(G[\M]\cong K_{k,l}\).
\end{corollary}

\begin{proof}
Behague proved this for all \((k,l)\ne(3,3)\) \cite[Theorem 2]{behague}. The
theorem above supplies the remaining case.
\end{proof}

\section{Verification method}

Although the construction is small enough to be checked directly from
Table~\ref{tab:coloring}, we describe the verification procedure explicitly.
For each color \(i\), form the map \(p_i\) sending a vertex \(v\) to the other
endpoint of the unique color-\(i\) edge incident with \(v\). The assertion
that \(M_i\) is a perfect matching is equivalent to saying that \(p_i\) is a
fixed-point-free involution on the 64 vertices. Since the table lists every
edge of \(Q_6\) exactly once, this also checks that the color classes partition
\(E(Q_6)\).

For a pair of colors \(i,j\), the graph \(M_i\cup M_j\) is then the graph
generated by the two involutions \(p_i\) and \(p_j\). Its connected components
are cycles. Starting from an unvisited vertex and alternating the two maps
therefore gives the length of one component; repeating this until all vertices
are visited gives the cycle-length multiset for the pair. Since \(Q_6\) has
64 vertices, the pair is Hamiltonian exactly when this multiset is \(\{64\}\).
The displayed cycle-length table in the proof is the complete output of this
procedure for the 15 unordered pairs of colors.

\section*{Verification files}

The table above is self-contained. For independent reproducibility, the
accompanying files are:
\begin{itemize}
\item \texttt{q6\_k33\_solution.json},
\item \texttt{verify\_q6\_k33.py},
\item \texttt{q6\_k33\_selfcheck.py}.
\end{itemize}
The first lists the six matchings explicitly as edge sets. The second checks
that the JSON file is a 1-factorization and computes all pairwise cycle
lengths. The third hard-codes Table~\ref{tab:coloring} and performs the same
check without reading any external certificate file. Thus the certificate and
the compact table can be checked independently.

\section*{Declaration of competing interest}

The author declares no known competing financial interests or personal
relationships that could have appeared to influence the work reported in this
paper.

\section*{Funding}

No funding was received for this work.

\section*{Data availability}

The construction and verification scripts are archived on Zenodo
(DOI: \href{https://doi.org/10.5281/zenodo.19881180}{10.5281/zenodo.19881180})
and are also available in
\href{https://github.com/dmzlxm/q6-k33-factorization}{the project repository}.

\section*{Declaration of generative AI and AI-assisted technologies in the writing process}

OpenAI Codex/ChatGPT was used to assist with computational search, generation
of verification scripts, and drafting/editing of the manuscript. The author
checked the construction and verification output and takes full responsibility
for the content of the paper. This statement should be adjusted to match the
target journal's required disclosure format.

\begin{thebibliography}{9}

\bibitem{anderson}
B.~A. Anderson,
Finite topologies and Hamiltonian paths,
Journal of Combinatorial Theory, Series B 14 (1973) 87--93.

\bibitem{behague}
N.~C. Behague,
Semi-perfect 1-factorizations of the hypercube,
Discrete Mathematics 342(6) (2019) 1696--1702.
\href{https://doi.org/10.1016/j.disc.2019.01.035}{doi:10.1016/j.disc.2019.01.035};
\href{https://arxiv.org/abs/1811.06389}{arXiv:1811.06389}.

\bibitem{bryant_maenhaut_wanless}
D.~E. Bryant, B.~Maenhaut, I.~M. Wanless,
A family of perfect factorisations of complete bipartite graphs,
Journal of Combinatorial Theory, Series A 98(2) (2002) 328--342.

\bibitem{chitra_muthusamy}
V.~Chitra, A.~Muthusamy,
A note on semi-perfect 1-factorization and Craft's conjecture,
Graph Theory Notes of New York 64 (2013) 22--25.

\bibitem{gochev_gotchev}
V.~S. Gochev, I.~S. Gotchev,
On \(k\)-semi-perfect one-factorizations of \(Q_n\) and Craft's conjecture,
Graph Theory Notes of New York 58 (2010) 36--41.

\bibitem{kotzig}
A.~Kotzig,
Hamilton graphs and Hamilton circuits,
in: Theory of Graphs and its Applications, Proceedings of the Symposium held
in Smolenice in June 1963, Czechoslovak Academy of Sciences, Prague, 1964,
pp. 63--82.

\bibitem{kralovic_kralovic}
R.~Kr\'alovi\v{c}, R.~Kr\'alovi\v{c},
On semi-perfect 1-factorizations,
in: Structural Information and Communication Complexity, Lecture Notes in
Computer Science 3499, Springer, Berlin, 2005, pp. 216--230.

\bibitem{laufer}
H.~B. Laufer,
On strongly Hamiltonian complete bipartite graphs,
Ars Combinatoria 9 (1980) 43--46.

\bibitem{stong}
R.~Stong,
Hamilton decompositions of directed cubes and products,
Discrete Mathematics 306(18) (2006) 2186--2204.

\end{thebibliography}

\end{document}
